\section{区間の反転}\label{節:区間の反転}

前節では,区間が原点について対称であった.では,原点を中心としない区間 $[a,b]$ でも,左右の値を組にして考えられないだろうか.本節では特に,区間の中央について対称な関数 $g$,すなわち
\[
g(a+b-x)=g(x)
\]
を満たす $g$ に一次因子 $x$ が掛かった積分
\[
\int_a^b xg(x)\dd x
\]
を考える.この積分から $x$ を取り除ければ,あとは $\int_a^b g(x)\dd x$ を計算するだけで済む.

少々発想しにくい変形であるため,結果からお見せしよう.
\begin{equation}\label{式:対称な重みの一次因子}
\int_a^b xg(x)\dd x
 =\frac{a+b}{2}\int_a^b g(x)\dd x
\end{equation}
が成り立つ.つまり,積分の中の $x$ を区間の中央 $\frac{a+b}{2}$ に置き換えてよいのだ.

これを示すために,区間 $[a,b]$ の両端を交換する置換 $x=a+b-t$ を考える.対応は
\[
x=a\longleftrightarrow t=b,\qquad
x=b\longleftrightarrow t=a
\]
であり,\,$\dd x=-\dd t$ だから,一般の関数 $f$ について
\[
\begin{aligned}
\int_a^b f(x)\dd x
 &=\int_b^a f(a+b-t)(-\dd t)\\
 &=\int_a^b f(a+b-t)\dd t
\end{aligned}
\]
となる.置換の直後は積分の上下端が逆になるが,微分の負号と合わせて向きを戻している.積分変数の文字を $x$ に戻せば,
\begin{equation}\label{式:区間反転}
\int_a^b f(x)\dd x=\int_a^b f(a+b-x)\dd x
\end{equation}
を得る.

$I=\int_a^b xg(x)\dd x$ と置き,式\eqref{式:区間反転}を $f(x)=xg(x)$ に使う.\,$g(a+b-x)=g(x)$ だから
\[
I=\int_a^b(a+b-x)g(x)\dd x
\]
である.もとの式と加えると,\,$x+(a+b-x)=a+b$ なので
\[
\begin{aligned}
2I&=\int_a^b\{x+(a+b-x)\}g(x)\dd x\\
 &=(a+b)\int_a^b g(x)\dd x
\end{aligned}
\]
となる.両辺を $2$ で割れば式\eqref{式:対称な重みの一次因子}を得る.

ここで行ったのは,反転した積分ともとの積分を加えて $2$ で割ることである.一般の $f$ について書けば
\begin{equation}\label{式:反転平均}
\int_a^b f(x)\dd x
 =\frac12\int_a^b\{f(x)+f(a+b-x)\}\dd x
\end{equation}
となる.式\eqref{式:反転平均}の中括弧が定数 $k$ になれば,答えは $\frac12k(b-a)$ である.もとの式を個別に積分できるかどうかは,この計算では問題にならない.次節以降では,中括弧が定数になる例を扱う.

\begin{bookpoint}
\textbf{POINT}\par
$x=a+b-t$ と置いて反転した積分を作り,もとの積分と加える.\,$g(a+b-x)=g(x)$ なら,\,$x$ を区間の中央 $\frac{a+b}{2}$ に置き換えられる.
\end{bookpoint}

\begin{example}{反転による計算}
次の定積分を求めよ.
\[
I=\int_0^\pi x\sin^3x\dd x.
\]
\end{example}

\paragraphbf{方針}
$\sin(\pi-x)=\sin x$ より,\,$\sin^3x$ は区間の中央 $\frac{\pi}{2}$ について対称である.反転で $x$ を $\frac{\pi}{2}$ に置き換え,残りは $u=\cos x$ と置く.

\paragraphbf{解答}
$I=\frac{2\pi}{3}$

\paragraphbf{解法}
$x$ と三角関数の積を見ると部分積分を考えたくなるだろう.しかしこの区間では $\sin(\pi-x)=\sin x$ であり,\,$g(x)=\sin^3x$ は $g(\pi-x)=g(x)$ を満たす.式\eqref{式:対称な重みの一次因子}で $a=0$,\,$b=\pi$ とすれば
\[
I=\frac{\pi}{2}\int_0^\pi\sin^3x\dd x
\]
となる.残りは $\sin^3x=(1-\cos^2x)\sin x$ として $u=\cos x$ と置く.\,$x=0,\pi$ に対して $u=1,-1$ であり,\,$\sin x\dd x=-\dd u$ だから
\[
\begin{aligned}
\int_0^\pi\sin^3x\dd x
 &=\int_1^{-1}(1-u^2)(-\dd u)\\
 &=\int_{-1}^1(1-u^2)\dd u\\
 &=\left[u-\frac{u^3}{3}\right]_{-1}^1\\
 &=\frac43
\end{aligned}
\]
である.したがって
\[
I=\frac{\pi}{2}\cdot\frac43=\frac{2\pi}{3}
\]
を得る.反転で $x$ を取り除き,置換で三角関数を取り除くという2段階で計算できた.

\paragraphbf{反転の発想}
以下,\,$x$ を $a+b-x$ に移す変換を $x\mapsto a+b-x$ のように矢印で書く.反転 $x\mapsto a+b-x$ は,前節の $x\mapsto-x$ を区間の中央 $c=\frac{a+b}{2}$ のまわりで行ったものである.実際,\,$x=c+s$ と置けば $a+b-x=2c-(c+s)=c-s$ であり,反転は $s\mapsto-s$ にほかならない.

この見方で式\eqref{式:対称な重みの一次因子}をもう一度導いてみよう.\,$h=\frac{b-a}{2}$ とし,\,$x=c+s$ と置くと,区間は $-h\le s\le h$ となる.\,$G(s)=g(c+s)$ と置けば,\,$g(a+b-x)=g(x)$ は $G(-s)=G(s)$,すなわち $G$ が偶関数であることを表す.したがって
\[
\begin{aligned}
\int_a^b xg(x)\dd x
 &=\int_{-h}^h(c+s)G(s)\dd s\\
 &=c\int_{-h}^hG(s)\dd s+\int_{-h}^h sG(s)\dd s
\end{aligned}
\]
となる.\,$sG(s)$ は奇関数なので第2項は $0$ であり,第1項は $c\int_a^b g(x)\dd x$ に等しい.中央からの距離が等しい2点では $g$ の値が等しく,その2点の $x$ 座標の平均は常に区間の中央 $c$ になる.式\eqref{式:対称な重みの一次因子}は,このことを積分で表したものである.

\paragraphbf{反転が使えない場合}
式\eqref{式:対称な重みの一次因子}を使う際には,\,$xg(x)$ 全体が対称なのではなく,\,$g$ が対称であることを確認する.たとえば $\int_0^1 x\e^x\dd x$ では $\e^{1-x}=\e^x$ が成り立たないので,\,$x$ を区間の中央 $\frac12$ に置き換えることはできない.

また,\,$x^2g(x)$ まで同じように中央の2乗に置き換えられるわけでもない.\,$x^2+(a+b-x)^2$ は $x$ を含むからである.反転して足すことで何が消えるのかを,その都度式で確かめる必要がある.
